\subsection{\texorpdfstring{第~\ref{sec:minimizing-properties-of-geodesics}~节}{测地线的最短性}}\label{subsec:solutions-minimizing-properties-of-geodesics}

\begin{description}[font=\normalfont,leftmargin=2em,labelwidth=1.5em]
\RGsolution{ex:bi-invariant-metrics-on-lie-groups}{bi-invariant-metrics-on-lie-groups}
\begin{enumerate}[label=(\alph*),ref=\alph*,leftmargin=*]
\item\label{sol:inversion-is-isometry-for-bi-invariant-metric} 设 $v\in\mathfrak{g}=T_eG$. 则
\begin{align*}
  (\mathrm{d}i)_ev
  &=\left.\frac{\mathrm{d}}{\mathrm{d}t}i(\exp(tv))\right|_{t=0}
   =\left.\frac{\mathrm{d}}{\mathrm{d}t}(\exp(tv))^{-1}\right|_{t=0}\\
  &=\left.\frac{\mathrm{d}}{\mathrm{d}t}(\exp(-tv))\right|_{t=0}=-v
\end{align*}
(其中 $\exp$ 为 Lie 群的指数映射), 故 $(\mathrm{d}i)_e=-\operatorname{id}$. 此外, 对 $g\in G$,
\[
  i(g)=g^{-1}=(h^{-1}g)^{-1}h^{-1}=(R_{h^{-1}}\circ i\circ L_{h^{-1}})(g),
\]
这说明对每个 $h\in G$, 有
\[
  (\mathrm{d}i)_g=(\mathrm{d}R_{h^{-1}})_{(h^{-1}g)^{-1}}(\mathrm{d}i)_{h^{-1}g}(\mathrm{d}L_{h^{-1}})_g.
\]
取 $h=g$, 得
\[
  (\mathrm{d}i)_g=(\mathrm{d}R_{g^{-1}})_e(\mathrm{d}i)_e(\mathrm{d}L_{g^{-1}})_g.
\]
因此, 对 $v,w\in T_gG$,
\begin{align*}
  &\langle(\mathrm{d}i)_gv,(\mathrm{d}i)_gw\rangle\\
  &\quad=\bigl\langle(\mathrm{d}R_{g^{-1}})_e(\mathrm{d}i)_e(\mathrm{d}L_{g^{-1}})_gv,
        (\mathrm{d}R_{g^{-1}})_e(\mathrm{d}i)_e(\mathrm{d}L_{g^{-1}})_gw\bigr\rangle\\
  &\quad=\langle(\mathrm{d}i)_e(\mathrm{d}L_{g^{-1}})_gv,(\mathrm{d}i)_e(\mathrm{d}L_{g^{-1}})_gw\rangle\\
  &\quad=\langle-(\mathrm{d}L_{g^{-1}})_gv,-(\mathrm{d}L_{g^{-1}})_gw\rangle=\langle v,w\rangle,
\end{align*}
这里用了 $R_{g^{-1}}$ 和 $L_{g^{-1}}$ 都是等距映射这一事实.
\item 设 $\widetilde c:=(c_v)^{-1}=i\circ c_v$. 由 (\ref{sol:inversion-is-isometry-for-bi-invariant-metric}) 得 $\widetilde c(0)=e$ 且 $\dot{\widetilde c}(0)=-v$, 因而 $(c_v)^{-1}(t)=c_{-v}(t)$ 是过 $e$ 且初速度为 $-v$ 的测地线. 另一方面, 若 $\gamma(t):=c_v(-t)$, 则 $\gamma(0)=c_v(0)=e$ 且 $\dot\gamma(0)=-v$, 故 $\gamma(t)=c_{-v}(t)$. 因此 $(c_v)^{-1}(t)=c_v(-t)$. 现设 $I:=(a,b)$ 是 $c_v$ 有定义的含 $0$ 的最大开区间, 并固定 $t_0\in I$ 使 $0<t_0<\varepsilon$, 其中 $\varepsilon>0$ 使 $\exp_e(B_\varepsilon(0))$ 为正规球. 则 $\gamma(t):=c_v(t_0)c_v(t)$ 也是定义在 $I$ 上的测地线 (因为 $L_{c_v(t_0)}$ 是等距映射), 且满足
\[
  \gamma(-t_0)=c_v(t_0)c_v(-t_0)=c_v(t_0)(c_v(t_0))^{-1}=e
\]
和 $\gamma(0)=c_v(t_0)$. 因为 $e$ 与 $c_v(t_0)$ 由两条长度均为 $t_0\lVert v\rVert<\varepsilon$ 的测地线连接, 所以这两条测地线必重合. 因而对所有使两边都有定义的 $t\in I$, 有 $\gamma(t)=c_v(t_0+t)$. 若 $b<+\infty$, 则 $\gamma$ 可将 $c_v$ 延拓到 $I$ 以外, 这是不可能的. 此外, 由 $c_v(-t)=(c_v(t))^{-1}$ 得 $a=-\infty$, 因而 $c_v$ 对所有 $t\in\mathbb{R}$ 都有定义. 因为对每对满足 $|t_0|<\varepsilon$ 的 $t,t_0\in\mathbb{R}$, 有
\[
  c_v(t_0+t)=\gamma(t)=c_v(t_0)c_v(t),
\]
给定 $s\in\mathbb{R}$ 后, 选取 $m>0$ 使 $\left|\frac{s}{m}\right|<\varepsilon$, 则有
\begin{align*}
  c_v(s+t)
  &=c_v\left(\frac{m}{m}s+t\right)=c_v\left(\frac{s}{m}\right)c_v\left(\frac{m-1}{m}s+t\right)\\
  &=\cdots=c_v\left(\frac{s}{m}\right)\cdots c_v\left(\frac{s}{m}\right)c_v(t)=c_v(s)c_v(t).
\end{align*}
\item 设 $v\in\mathfrak{g}$, 考虑满足 $c_v(0)=e$ 的测地线 $c_v:I\to G$. 设 $X$ 为满足 $X_e=v$ 的左不变向量场. 因为 $c_v(s+t)=c_v(s)c_v(t)$, 所以
\begin{align*}
  \dot c_v(s)&=\left.\frac{\mathrm{d}}{\mathrm{d}t}c_v(s+t)\right|_{t=0}
    =\left.\frac{\mathrm{d}}{\mathrm{d}t}(c_v(s)c_v(t))\right|_{t=0}\\
  &=(\mathrm{d}L_{c_v(s)})_{c_v(0)}\dot c_v(0)
    =(\mathrm{d}L_{c_v(s)})_ev=X_{c_v(s)}.
\end{align*}
因此 $c_v$ 是 $X$ 的满足 $c_v(0)=e$ 的积分曲线. 另一方面, 因为 $L_g$ 是等距映射, $G$ 的所有测地线都是过 $e$ 的测地线的左平移像. 又因为
\[
  (L_g\circ c_v)(0)=gc_v(0)=g
\]
和
\[
  \left.\frac{\mathrm{d}}{\mathrm{d}t}(L_g\circ c_v)(t)\right|_{t=0}
  =(\mathrm{d}L_g)_e\dot c_v(0)=(\mathrm{d}L_g)_ev=X_g,
\]
所以所有测地线都是左不变向量场的积分曲线. 最后,
\[
  \exp_e(v)=c_v(1)=\exp(v).
\]
\item 设 $X$ 是左不变向量场. 因为它的积分曲线是测地线, 所以 $\nabla_XX=0$. 因此, 若 $X$ 和 $Y$ 都是左不变向量场, 则
\[
  0=\nabla_{X+Y}(X+Y)=\nabla_XY+\nabla_YX=2\nabla_XY-[X,Y],
\]
这里使用了 Levi–Civita 联络的对称性.
\item 为验证 $S^3\cong SU(2)$ 中的标准度量也右不变, 注意若 $v\in SU(2)$ 是单位四元数, 即 $v^*v=1$, 则
\[
  \langle q_1v,q_2v\rangle=\operatorname{Re}(q_1vv^*q_2^*)=\operatorname{Re}(q_1q_2^*)=\langle q_1,q_2\rangle.
\]
为验证 $\mathcal{M}_{n\times n}\cong\mathbb{R}^{n^2}$ 的 Euclid 度量在 $O(n)$ 上诱导的度量也右不变, 注意若 $S\in O(n)$, 则
\[
  \langle AS,BS\rangle=\operatorname{tr}(ASS^tB^t)=\operatorname{tr}(AB^t)=\langle A,B\rangle.
\]
\item 设 $G$ 是 $n$ 维紧 Lie 群, $\langle\cdot,\cdot\rangle$ 是 $G$ 上的左不变度量. 给定 $v,w\in T_gG$, 定义
\[
  \langle\langle v,w\rangle\rangle=\int_G f,
\]
其中积分关于 (左不变的) Riemann 体积元取得, 而函数 $f:G\to\mathbb{R}$ 为
\[
  f(h)=\langle(\mathrm{d}R_h)_gv,(\mathrm{d}R_h)_gw\rangle.
\]
容易验证 $\langle\langle\cdot,\cdot\rangle\rangle$ 定义了 $G$ 上的 Riemann 度量 (其正性来自当 $v=w\ne0$ 时 $f>0$). 另一方面, 它左不变, 因为
\begin{align*}
  &\langle(\mathrm{d}R_h)_{kg}(\mathrm{d}L_k)_gv,(\mathrm{d}R_h)_{kg}(\mathrm{d}L_k)_gw\rangle\\
  &\quad=\langle(\mathrm{d}(R_h\circ L_k))_gv,(\mathrm{d}(R_h\circ L_k))_gw\rangle\\
  &\quad=\langle(\mathrm{d}(L_k\circ R_h))_gv,(\mathrm{d}(L_k\circ R_h))_gw\rangle\\
  &\quad=\langle(\mathrm{d}L_k)_{gh}(\mathrm{d}R_h)_gv,(\mathrm{d}L_k)_{gh}(\mathrm{d}R_h)_gw\rangle\\
  &\quad=\langle(\mathrm{d}R_h)_gv,(\mathrm{d}R_h)_gw\rangle,
\end{align*}
所以积分后得
\[
  \langle\langle(\mathrm{d}L_k)_gv,(\mathrm{d}L_k)_gw\rangle\rangle=\langle\langle v,w\rangle\rangle.
\]
最后, 它也右不变, 因为
\begin{align*}
  &\langle(\mathrm{d}R_h)_{gk}(\mathrm{d}R_k)_gv,(\mathrm{d}R_h)_{gk}(\mathrm{d}R_k)_gw\rangle\\
  &\quad=\langle(\mathrm{d}(R_h\circ R_k))_gv,(\mathrm{d}(R_h\circ R_k))_gw\rangle\\
  &\quad=\langle(\mathrm{d}R_{kh})_gv,(\mathrm{d}R_{kh})_gw\rangle=f(kh),
\end{align*}
从而
\[
  \langle\langle(\mathrm{d}R_k)_gv,(\mathrm{d}R_k)_gw\rangle\rangle
  =\int_G f\circ L_k=\int_G f=\langle\langle v,w\rangle\rangle
\]
(因为体积元左不变).
\end{enumerate}

\RGsolution{ex:jacobi-fields-and-cut-locus}{jacobi-fields-and-cut-locus}
\begin{enumerate}[label=(\alph*),ref=\alph*,leftmargin=*]
\item\label{sol:jacobi-identity-for-geodesic-coordinates} 有
\[
  R(X,Y_i)X=\nabla_X\nabla_{Y_i}X-\nabla_{Y_i}\nabla_XX-\nabla_{[X,Y_i]}X.
\]
因为 $X$ 和 $Y_i$ 满足 $\nabla_XX=0$ 和 $[X,Y_i]=0$, 所以
\[
  R(X,Y_i)X=\nabla_X\nabla_{Y_i}X=\nabla_X([Y_i,X]+\nabla_XY_i)=\nabla_X\nabla_XY_i.
\]
\item 设
\[
  Y(t):=\left.\frac{\partial}{\partial s}\exp_p(tv(s))\right|_{s=0},
\]
其中 $v:(-\varepsilon,\varepsilon)\to T_pM$ 满足 $v(0)=\dot c(0)$. 由于
\[
  \exp_p(tv(0))=\exp_p(t\dot c(0))=c(t),
\]
可知 $Y$ 是沿 $c$ 的向量场. 此外, $\exp_p(0v(s))=p$, 蕴含 $Y(0)=0$. 若 $\{\dot c(t_0),Y(t_0)\}$ 线性无关, 则可将 $(t,s)$ 扩充为 $c(t_0)$ 邻域中的一组局部坐标, 并利用 (\ref{sol:jacobi-identity-for-geodesic-coordinates}) 证明 $Y$ 在 $t_0$ 的某个邻域内满足沿 $c$ 的 Jacobi 方程. 现在证明, 使 $\{\dot c(t_0),Y(t_0)\}$ 线性相关的点 $t_0$ 的集合 $D\subset\mathbb{R}$ 或者内部为空 (于是由连续性, $Y$ 满足沿 $c$ 的 Jacobi 方程), 或者对某个 $\beta\in\mathbb{R}$, 有 $Y(t)=\beta t\dot c(t)$ (这也是 Jacobi 方程的一个解, 因为 $R(X,X)X=0$).

的确, 假设 $D$ 的内部非空, 并设 $(a,b)$ 是其内部的一个连通分支. 此时对某个函数 $f:(a,b)\to\mathbb{R}$ 有 $Y(t)=f(t)\dot c(t)$, 于是 $\frac{DY}{\mathrm{d}t}(t)=\dot f(t)\dot c(t)$. 因为 $D$ 闭, 若 $b$ 有限, 则 $Y$ 会在以 $b$ 为下确界的某个开区间上是沿 $c$ 的 Jacobi 方程的一个解, 满足 $Y(b)=\alpha\dot c(b)$ 和 $\frac{DY}{\mathrm{d}t}(b)=\beta\dot c(b)$ (其中 $\alpha:=f(b)$, $\beta:=\dot f(b)$). 由于
\[
  \widetilde Y(t):=(\alpha+\beta(t-b))\dot c(t)
\]
也是一个解, 且满足 $\widetilde Y(b)=\alpha\dot c(b)$ 及 $\frac{D\widetilde Y}{\mathrm{d}t}(b)=\beta\dot c(b)$, Picard–Lindelöf 定理蕴含在该区间上 $Y(t)=\widetilde Y(t)$, 从而 $b$ 将不是该连通分支的上确界 (因为 $\widetilde Y(t)$ 和 $\dot c(t)$ 线性相关). 于是 $b=+\infty$. 类似地, 有 $a=-\infty$, 即 $D=\mathbb{R}$. 现在 $Y(t)$ 的公式可写成
\[
  Y(t)=(\mathrm{d}\exp_p)_{t\dot c(0)}(t\dot v(0)),
\]
所以
\[
  \dot v(0)=\lim_{t\to0}\frac1tY(t)
    =\lim_{t\to0}\frac{f(t)}{t}\dot c(t)=\beta\dot c(0)=\beta v(0),
\]
其中
\[
  \beta=\lim_{t\to0}\frac{f(t)}{t}.
\]
还应注意, $Y(t)$ 只依赖于 $v(0)$ 和 $\dot v(0)$. 因此可以选取例如 $v(s)=(1+s\beta)v(0)$, 此时
\[
  \exp_p(tv(s))=\exp_p((1+s\beta)tv(0))=c((1+s\beta)t),
\]
故 $Y(t)=\beta t\dot c(t)$.

反过来, 给定满足 $Y(0)=0$ 的沿 $c$ 的 Jacobi 方程的解 $Y(t)$, 选取
\[
  \widetilde Y(t)=\left.\frac{\partial}{\partial s}\exp_p(t(v_0+sw_0))\right|_{s=0},
\]
其中 $v_0=\dot c(0)$, $w_0=\frac{DY}{\mathrm{d}t}(0)$. 则 $\widetilde Y(0)=0$. 此外, 利用正规坐标的 Christoffel 符号在原点处为零, 容易看出
\[
  \frac{D\widetilde Y}{\mathrm{d}t}(0)=w_0=\frac{DY}{\mathrm{d}t}(0).
\]
如上所见, $\widetilde Y(t)$ 也是沿 $t$ 的 Jacobi 方程的一个解, 由 Picard–Lindelöf 定理, 对所有 $t\in\mathbb{R}$, 有 $\widetilde Y(t)=Y(t)$.
\item 这直接由 Jacobi 场的公式
\[
  Y(t)=(\mathrm{d}\exp_p)_{t\dot c(0)}(t\dot v(0))
\]
得出.
\item 由 Jacobi 方程得
\begin{align*}
  \frac{\mathrm{d}^2}{\mathrm{d}t^2}\bigl(\lVert Y(t)\rVert^2\bigr)
  &=2\frac{\mathrm{d}}{\mathrm{d}t}\left\langle\frac{DY}{\mathrm{d}t}(t),Y(t)\right\rangle\\
  &=2\left\langle\frac{D^2Y}{\mathrm{d}t^2}(t),Y(t)\right\rangle
      +2\left\langle\frac{DY}{\mathrm{d}t}(t),\frac{DY}{\mathrm{d}t}(t)\right\rangle\\
  &=2\langle R(\dot c(t),Y(t))\dot c(t),Y(t)\rangle
      +2\left\lVert\frac{DY}{\mathrm{d}t}(t)\right\rVert^2\ge0.
\end{align*}
此外, 因为 $\frac{DY}{\mathrm{d}t}(0)\ne0$ (否则 $Y\equiv0$), 此导数在 $t=0$ 的某邻域内显然严格为正. 所以 $Y(t)$ 只能有一个零点 (在 $t=0$ 时).
\item
\begin{enumerate}[label=(\roman*),ref=\roman*,leftmargin=*]
\item 把 $T^n$ 看成在 $(n-1)$ 维面上作通常等同的 $[0,1]^n$, 则点 $p=(\frac12,\ldots,\frac12)$ 的割迹恰为各 $(n-1)$ 维面之并. 从 $p$ 到这些面上的每个点, 至少有两条长度相同而像不同的测地线.
\item $p\in S^n$ 的割迹由对径点 $-p$ 构成, 它显然与 $p$ 共轭.
\item 把 $\mathbb{R}P^n$ 看成将赤道上的对径点等同后的 $S^n$ 北半球, 则北极的割迹恰为赤道. 从北极到赤道上的每一点, 恰有两条长度相同而像不同的测地线.
\end{enumerate}
\end{enumerate}
\end{description}
